\section{Step 36: Publish index metadata}
Store admitted count, period, duration, complete file bytes and directory
entry count at context offsets 24 through 56. Return zero and eight times the
frame count. Retain directory capacity, scratch pointer, scratch capacity and
the sixteen-byte suffix. The supplied epilogue restores saved state and returns
to the caller, which reports that the directory is ready.

Fill \texttt{index\_publish}. Run the four valid fixtures and all failure
controls. Then reuse the same output storage for success, a rejected late command,
and another valid file. Rejection must retain the earlier directory and metadata;
a shorter later directory must retain bytes after its own active entries. The
scratch image has no such retention promise. A successful index now gives the
future seeker validated record locations and timing metadata.

\paragraph{Observe} Prepare \texttt{s36-1} and stop before
\texttt{index\_published}. At \texttt{index\_published}, the return values are zero and
64. Published values are eight frames, period four, duration forty, file count
24\,795 and eight directory entries. The eight offsets are 24, 12\,328,
12\,355, 12\,395, 12\,415, 24\,719, 24\,752 and 24\,775.

\Needspace{7\baselineskip}
\paragraph{Check} What must a failed reindex retain?
\begin{itemize}
\item[A.] Only its scratch image.
\item[B.] The earlier public directory and context.
\item[C.] Only the input count.
\end{itemize}
